\documentclass{resume} % Use the custom resume.cls style

\usepackage[left=0.4 in,top=0.4in,right=0.4 in,bottom=0.4in]{geometry} % Document margins
\newcommand{\tab}[1]{\hspace{.2667\textwidth}\rlap{#1}} 
\newcommand{\itab}[1]{\hspace{0em}\rlap{#1}}
\name{Xinchang Lin} % Your name
% You can merge both of these into a single line, if you do not have a website.
\address{+31 6 84656778 \\ Eindhoven. NL}
\address{\href{mailto:xlin@mailbox.org}{xlin@mailbox.org}  \\  \href{mailto:x.lin2@student.tue.nl}{x.lin2@student.tue.nl} \\ \href{https://www.linkedin.com/in/xinchang-lin/}{linkedin.com/in/xinchang-lin}  }  %

\begin{document}

%----------------------------------------------------------------------------------------
%	OBJECTIVE
%----------------------------------------------------------------------------------------

%\begin{rSection}{OBJECTIVE}

%{Software Engineer with 2+ years of experience in XXX, seeking full-time XXX roles.}


%\end{rSection}
%----------------------------------------------------------------------------------------
%	EDUCATION SECTION
%----------------------------------------------------------------------------------------

\begin{rSection}{Education}

{\bf BSc Mechanical Engineering}, Technische Universiteit Eindhoven \hfill {2025 - Expected 2028}\\
Relevant Coursework: Mechanics, Dynamics, Calculus, Introduction transport phenomena, Structures and properties of materials, Principles of design and programming\\
Average: 8.0/10

{\bf Psychology}, South China Normal University \hfill {2022 - 2025}\\
Average: 88.5/100
%Minor in Linguistics \smallskip \\
%Member of Eta Kappa Nu \\
%Member of Upsilon Pi Epsilon \\


\end{rSection}

%----------------------------------------------------------------------------------------
% TECHINICAL STRENGTHS	
%----------------------------------------------------------------------------------------
\begin{rSection}{SKILLS}

\begin{tabular}{ @{} >{\bfseries}l @{\hspace{6ex}} l }
Technical Skills & Python, Matlab, Marc Mentat, Siemens NX, Arduino, IBM SPSS
\\
Soft Skills & Communication, project coordination, workflow improvement, logistic management\\
Language & Chinese(native), English(C1), Dutch(elementary), Japanese(elementary)\\
\end{tabular}\\
\end{rSection}


%----------------------------------------------------------------------------------------
%	WORK EXPERIENCE SECTION
%---------------------------------------------------------------------------------------

 




\begin{rSection}{PROJECTS}
\vspace{-1.25em}
\item \textbf{Design of a launching mechanism.} {Built a simple launching mechanism with limited energy source, using trajectory analysis with ODE in Python and Matlab and CAD design for components in Siemens NX.}
\item \textbf{Multiped robot design.} {Built a robot that finishes a designated route, using inverse kinematics modelling in Siemens Simcenter Motion, CAD design for components and Arduino programming of sensor and servo motor control.}
\item \textbf{Design of a truss structure.} {Built a truss structure to sustain certain amount of load, using static FEA, calculations and simulations of force distribution in Mentat.}
\end{rSection} 

%----------------------------------------------------------------------------------------
\begin{rSection}{Extra-Curricular Activities} 
\begin{itemize}
    \item 	Creating \href{https://obsp.de/}{blog posts} on technical issues, gaining 4K unique visitors per month.
    \item	Community webhosting, providing several instances of free and open source software to a group of people, with skills in coordinating user needs, resource allocation, and basic service maintenance. Implemented minor features and troubleshooting improvements for hosted services.
\end{itemize}


\end{rSection}

%----------------------------------------------------------------------------------------
\begin{rSection}{EXPERIENCE}

\textbf{Private tutor} \hfill Sept 2024 - Mar 2025\\
Self-employed \hfill \textit{Shantou, CN}
 \begin{itemize}
    \itemsep -3pt {} 
     \item Tutored primary school students in mathematics and science.
     \item Expanded students’ knowledge and encouraged curiosity about advanced problems.
 \end{itemize}
 \textbf{Secretary-General’s Office Staff} \hfill Sept 2022 - Jun 2023\\
Psychological Research Society, School of Psychology of SCNU \hfill \textit{Guangzhou, CN}
 \begin{itemize}
    \itemsep -3pt {} 
     \item Responsible for the logistics management of the society
    \item Tracked inventory and expenses
    \item Improving workflow between different departments
 \end{itemize}
\end{rSection} 


\end{document}
